% 原書374--378頁、light-cone gauge から path integral まで。
\gloss{光円錐ゲージ（light-cone gauge）}
{弦理論では、世界面の時間座標を特定の時空のヌル座標と一致させる選び方。局所対称性を持つ理論では、与えられたヌル方向の接続がゼロとなるゲージの選び方。}

\gloss{線形ディラトン理論（linear dilaton theory）}
{エネルギー運動量テンソルにスカラーの2階微分に比例する項を含む、スカラーCFT。弦理論では、ディラトンが位置の線形関数である場合に現れる。}

\gloss{線形シグマ模型（linear sigma model）}
{運動項が場に依存しないが、長距離の物理が非線形シグマ模型によって支配されるスカラー場理論。}

\gloss{Liouville場理論（Liouville field theory）}
{指数関数的相互作用を持つスカラー場のCFT。非臨界弦など、さまざまな状況で現れる。線形ディラトンと指数関数的タキオンの背景におけるボソン弦理論に対応する。}

\gloss{リトル弦理論（little string theory）}
{重力を含まず、相互作用するいくつかの弦理論のいずれか。特に、弦結合をゼロにする極限のNS5ブレーン上に見いだされる。}

\gloss{ループ展開（loop expansion）}
{量子場理論ではFeynman図による展開であり、$\hbar$のべきによる展開に等しい。\emph{弦のループ展開}は、無次元の弦結合$g$を用いたRiemann面についての和である。\emph{世界面のループ展開}は非線形シグマ模型の摂動展開であり、コンパクト化半径を$R_{\mathrm c}$として、$\alpha'/R_{\mathrm c}^2$のべきによる展開である。}

\gloss{Lorentz型（Lorentzian）}
{$(-,\ldots,-,+,\ldots,+)$という混合した符号を持つこと。}

\gloss{下降演算子（lowering operator）}
{与えられた状態のエネルギーを下げる演算子。CFTでは、Virasoro生成子$L_0$または$\widetilde L_0$を$n$だけ下げる演算子は、次数、すなわち下付き添字$n$を持つ。}

\gloss{M理論（M-theory）}
{1．狭義には、11次元のPoincaré不変性を持つ、IIA理論の強結合極限。2．より一般的な広義の用法では、さまざまな弱結合の弦理論と狭義のM理論を極限として含む、量子論全体。この名前は、理論の性質が未知であることを反映して、意図的に曖昧にされている。Mの意味としては、膜（membrane）、行列（matrix）、母（mother）、謎（mystery）などが提案されている。}

\gloss{M2ブレーン（M2-brane）}
{11次元超重力のポテンシャル$A_3$と結合する、M理論の2ブレーン。}

\gloss{M5ブレーン（M5-brane）}
{11次元超重力のポテンシャル$A_3$の磁荷を持つ、M理論の5ブレーン。}

\gloss{巨視的な弦（macroscopic string）}
{長さが弦の特徴的な長さスケールよりはるかに大きい弦。特に、時空を横切って延びる無限に長い弦を考えると便利なことがある。}

\gloss{Majorana条件（Majorana condition）}
{スピノル場に課すLorentz不変な実条件。時空次元が8を法として1、2、3、4または8の場合にのみ課すことができる。}

\gloss{周辺相互作用（marginal interaction）}
{無次元の強さが長さの尺度に依存しない相互作用。一般には、これは相互作用の結合の1次でだけ成立する場合がある。真に周辺的な相互作用とは、有限の結合でも周辺的であり続ける相互作用である。}

\gloss{行列模型（matrix models）}
{行列の自由度を持つ量子力学系。その臨界点は非臨界弦理論によって支配される。}

\gloss{行列理論（matrix theory）}
{$d=10$の超対称$U(n)$ Yang--Mills理論の次元縮約から得られる、行列の自由度と32個の超電荷を持つ量子力学系。大きな$n$の極限では、広義のM理論を定義すると予想されている。}

\gloss{極小模型（minimal models）}
{すべての共形族が縮退している、可解なCFTのいくつかの族。$N=0$、$N=1$、$N=2$の超共形対称性を持つユニタリーな極小模型には無限系列があり、それぞれ中心電荷1、$\frac32$、3に下から収束する。}

\gloss{Minkowski型（Minkowskian）}
{$(-,+,+,\ldots,+)$という符号を持つこと。}

\gloss{ミラー対称性（mirror symmetry）}
{異なる多様体上にコンパクト化した弦理論の同等性。この同等性は、一方の$N=2$超共形代数の$U(1)$電荷の符号を反転させるので、多様体のEuler数の符号を変える。}

\gloss{Möbius群（Möbius group）}
{球面上で大域的に定義された$SL(2,\mathbf C)$共形対称性、または円板上で大域的に定義された$SL(2,\mathbf R)$共形対称性。}

\gloss{モード演算子（mode operators）}
{量子場の空間的Fourier成分。}

\gloss{模型（model）}
{真空を参照。}

\gloss{モジュラー群（modular group）}
{大きな座標変換の群。しばしばRiemann面に適用されるが、時空にも適用できる。}

\gloss{モジュラー不変性（modular invariance）}
{大きな座標変換の下での、弦の経路積分の不変性。}

\gloss{モジュライ（moduli）}
{1．多様体の幾何を指定するパラメータ。代表的な例は、弦の世界面の複素構造のパラメータと、コンパクト化の幾何のパラメータである。2．量子場理論の、縮退した、通常は物理的に非同等な真空の空間を指定するパラメータ。これはコンパクト化の例と密接に関係する。古典極限のまわりで展開すると、場の方程式の各コンパクトな解が、主要次数で量子論の真空を与える。3．これらのパラメータの位置依存性に対応する、質量のない場。Goldstoneボソンと対比されたい。}

\gloss{モジュライ空間（moduli space）}
{モジュライを座標とする、幾何または真空の空間。}

\gloss{モノドロミー（monodromy）}
{局所的には1価な量について、非自明な閉経路のまわりで生じる多価性。}

\gloss{Montonen--Olive双対性（Montonen--Olive duality）}
{$d=4,\ N=4$ Yang--Mills理論の弱結合・強結合双対性。}

\gloss{Nambu--Goto作用（Nambu--Goto action）}
{時空内の世界面の不変面積に比例する弦の作用。}

\gloss{Narainコンパクト化（Narain compactification）}
{左向き・右向きの運動量の格子による、トーラス・コンパクト化の抽象的記述。}

\gloss{自然さの問題（naturalness problem）}
{自然定数が、ゼロでないと見積もられる寄与よりもはるかに小さい値を取る理由を説明する問題。例として、Higgsスカラーの質量、宇宙定数、QCDの$\theta$角がある。}

\gloss{Neumann境界条件（Neumann boundary condition）}
{境界で場の法線微分をゼロとする条件。場の値は自由に揺らぐ。}

\gloss{Neveu--Schwarz代数（Neveu--Schwarz algebra）}
{超カレントが反周期的であり、そのモードが半整数値となるセクターで、超カレントとエネルギー運動量テンソルのFourierモードがなす世界面の代数。}

\gloss{Neveu--Schwarz境界条件、NS境界条件（Neveu--Schwarz boundary condition）}
{閉弦、または開弦を二重化したものの世界面で、フェルミオン場が反周期的であるという条件。二重化の方法を参照。このときFourierモードは半整数値となる。}

\gloss{Neveu--Schwarz 5ブレーン（Neveu--Schwarz 5-brane）}
{I型およびII型超弦理論において、NS--NSの2形式ポテンシャルの磁荷を持つ5ブレーン。}

\gloss{Neveu--Schwarz--Neveu--Schwarz状態、NS--NS状態（Neveu--Schwarz--Neveu--Schwarz states）}
{II型超弦理論における、左向き・右向きの部分がともにボソン的である、ボソン的な閉弦状態。重力子とディラトンが含まれ、II型の場合には2形式ポテンシャルも含まれる。}

\gloss{ノーゴースト定理（no-ghost theorem）}
{1．OCQまたはBRSTのHilbert空間が正の内積を持つという定理。2．さらに、弦の振幅がこの空間で良く定義され、ユニタリーであるという定理。}

\gloss{ノースケール模型（no-scale model）}
{ある近似の範囲で、超対称性が破れた縮退真空の線を持つ場理論。}

\gloss{Noetherの定理（Noether's theorem）}
{Lagrangianの不変性から保存量が導かれるという定理。}

\gloss{非可換幾何（noncommutative geometry）}
{空間上の関数の代数に着目した、通常の幾何の一般化。Dブレーンの非可換な集団座標は、このような一般化の必要性を示唆する。}

\gloss{非臨界弦理論（noncritical string theory）}
{1．Weyl不変でない弦理論、すなわち、時空への埋め込みとは独立な世界面上の距離の測り方を持つ弦理論。独立な世界面計量場を持つ弦や、短距離カットオフを持つ弦が含まれる。2．より最近では、Weyl不変な世界面理論を伴う弱結合極限を持たない任意の弦理論にも、この用語が使われている。この意味では、$(2,0)$理論やリトル弦理論のように、結合が1程度に固定された、弦状の励起を持つさまざまな理論を含む。そのような理論では、世界面の位相を変える過程を止められないため、世界面は明確には定義されない。}

\gloss{非線形シグマ模型（nonlinear sigma model）}
{運動項の係数が場に依存するスカラー場理論。曲がった場の空間による自然な解釈を持ち、曲がった時空、またはより一般には位置に依存する背景場の中の弦理論に対応する。}

\gloss{非繰り込み定理（nonrenormalization theorem）}
{与えられた振幅または有効作用への量子補正の形を制限する定理。補正が消えること、摂動論の特定の次数でしか現れないこと、あるいは非摂動的にしか現れないことを要求する場合がある。}

\gloss{正規順序（normal ordering）}
{発散項を特定の仕方で差し引くことによって自由場の積を定義する処方。$:\ :$で表す\emph{共形正規順序}は、単純な共形変換性を持つ演算子を作る。$\canorm{\ }$で表す\emph{生成消滅正規順序}は、下降演算子を上昇演算子の右に置き、行列要素が単純な演算子を作る。$\bndnorm{\ }$で表す\emph{境界正規順序}は、共形正規順序に鏡像電荷の差引きを加えたもので、境界に近づいても有限な演算子を作る。}

\gloss{ヌル状態（null state）}
{自分自身を含む、すべての物理状態と直交する物理状態。または、共形族のすべての状態と直交する、その共形族の子孫状態。}

\gloss{旧共変量子化、OCQ（old covariant quantization）}
{電磁力学のGupta--Bleuler量子化に似た弦理論の量子化法。光円錐量子化およびBRST量子化と等価である。}

\gloss{1ループ（one-loop）}
{主要な量子補正。弦の摂動論ではEuler数がゼロの曲面から生じる。}

\gloss{開弦（open string）}
{位相的に線分である弦。}

\gloss{演算子方程式（operator equation）}
{量子論において、任意の行列要素の中で成立する演算子間の等式。同じことだが、任意の境界条件を持つ汎関数積分へ挿入したときに成立する等式。}

\gloss{演算子積展開、OPE（operator product expansion）}
{演算子の積を局所演算子の和として展開すること。その積を含む任意の期待値に対し、演算子間の距離がゼロになるときの漸近展開を与える。CFTでは展開は収束する。}

\gloss{オービフォールド（orbifold）}
{1．名詞として、多様体$M$の離散対称性の群を$H$とした商空間$M/H$。この商空間は$H$の固定点で特異となる。2．名詞として、世界面の離散対称群$H$をゲージ化して得られるCFTまたは弦理論。$H$の元が時空対称性なら、結果は商空間$M/H$上を伝播する弦の理論である。\emph{非可換オービフォールド}は、点群が非可換なものをいう。\emph{非対称オービフォールド}は、$H$が時空的解釈を持たず、一般に弦の右向きと左向きのモードに異なる作用をするものである。3．動詞として、$H$をゲージ化してこのようなCFTまたは弦理論を作ること。「捩り」の第3の定義と同義。}

\gloss{向き付けられた弦理論（oriented string theory）}
{世界面が定まった向きを持つ弦理論。世界面のパリティ反転をゲージ対称性として扱わない。}

\gloss{オリエンティフォールド（orientifold）}
{世界面の対称群$H$をゲージ化して作る弦理論。ただし$H$は、世界面のパリティ反転$\Omega$と他の対称性を組み合わせた元を含む。}

\gloss{オリエンティフォールド面（orientifold plane）}
{オリエンティフォールド群$H$、より具体的には$\Omega$を含む$H$の元の固定点からなる、任意の次元$p$の平面。}

\gloss{$p$ブレーン（p-brane）}
{空間的に$p$次元に広がった対象。例として、ブラック$p$ブレーン、D$p$ブレーン、M2およびM5ブレーン、NS5ブレーン、ヘテロティック5ブレーン、そして$d>4$ではKaluza--Kleinモノポールがある。}

\gloss{$p$形式（p-form）}
{完全反対称な$p$個の添字を持つテンソル。通常、添字を使わない記法で書く。}

\gloss{$p$形式ゲージ場（p-form gauge field）}
{$p$形式ポテンシャル$A$、$(p-1)$形式ゲージパラメータ$\lambda$、$(p+1)$形式の場の強さ$F$を持つ、可換ゲージ理論の一般化。$p=0$では通常の質量のないスカラー、$p=1$では可換ゲージ場である。$(p+1)$形式ポテンシャルは、世界体積上での形式の積分を通じて、$p$ブレーンと自然に結合する。\emph{自己双対$p$形式理論}とは$*F=F$である理論である。このためには時空次元が$d=2p+2$でなければならず、さらにMinkowski型の場合、$**=1$となるには$d$が4を法として2でなければならない。}

\gloss{$p$--$p'$弦（p--p' string）}
{一方の端点がD$p$ブレーン、他方の端点がD$p'$ブレーン上にある開弦。}

\gloss{$(p,q)$形式（(p,q)-form）}
{複素多様体上で、$p$個の正則添字と$q$個の反正則添字についてそれぞれ完全反対称なテンソル。}

\gloss{$(p,q)$弦（(p,q) string）}
{IIB理論における、$p$本のF弦と$q$本のD弦の束縛状態。}

\gloss{パラフェルミオンCFT（parafermion CFT）}
{$\mathbf Z_2$不変な自由フェルミオン理論を一般化する、$\mathbf Z_n$対称性を持つコセットCFTの族。$\mathbf Z_n$対称性を持つ系の一般的な臨界挙動を記述する。}

\gloss{パリティ変換（parity transformation）}
{1つ、または任意の奇数個の空間次元を反転する操作。時空パリティ$P$または$\beta$と、世界面パリティ$\Omega$を区別する。}

\gloss{分配関数（partition function）}
{量子系のスペクトルについて、Hamiltonianを$H$、温度を$T$として、$e^{-H/T}$を重みとした和。しばしば他の電荷も指数に含める。平衡統計力学の基本的な量である。弦理論ではトーラスまたは円筒上の経路積分で与えられ、したがって1ループ振幅に現れる。}

\gloss{経路積分（path integral）}
{量子系の遷移振幅を、あり得るすべての履歴についてのコヒーレントな和として表す方法。量子力学では履歴は粒子の経路、量子場理論では場の空間の経路、第一量子化された弦理論では弦の世界面の時空への埋め込みである。}
